% ============================================================
% Methodology. A reviewer should be able to re-implement the
% core of the method from this section alone. Structure:
% overview (with figure) → formalism → algorithm → details.
% Every design choice gets a WHY, not just a what.
% ============================================================
\section{Methodology}
\label{sec:methodology}

\subsection{Overview}
\label{subsec:overview}

% 1–2 paragraphs walking through the whole approach at a high
% level, ideally referencing an architecture figure. Describe
% the pipeline stages and how data flows between them:
%
% \begin{figure}[htbp]
%     \centering
%     \includegraphics[width=\linewidth]{figures/overview.pdf}
%     \caption{Overview of the proposed approach. ...}
%     \label{fig:overview}
% \end{figure}
%
% End with a roadmap: "Section~\ref{subsec:formulation} formalizes
% the problem; Section~\ref{subsec:algorithm} details the
% algorithm; Section~\ref{subsec:implementation} discusses
% implementation."
Give the reader the shape of the approach before the details.

\subsection{Problem Formulation}
\label{subsec:formulation}

% Formalize inputs, outputs, and objective. Define notation once
% here (or in a notation table) and use it consistently:
% \begin{equation}
%     \label{eq:objective}
%     \hat{\theta} = \arg\min_{\theta} \; \mathcal{L}(\theta; \mathcal{D})
% \end{equation}

\subsection{Proposed Algorithm}
\label{subsec:algorithm}

% The core contribution. Present the algorithm in pseudocode
% when it has non-trivial control flow:
%
% \begin{algorithm}[htbp]
%     \caption{Proposed Method}
%     \label{alg:proposed}
%     \begin{algorithmic}[1]
%         \Require input data $\mathcal{D}$, parameters $\lambda$
%         \Ensure trained model $M$
%         \State initialize $\theta$ randomly
%         \For{$t = 1$ to $T$}
%             \State sample mini-batch $B \subset \mathcal{D}$
%             \State update $\theta \gets \theta - \eta \nabla \mathcal{L}_B(\theta)$
%         \EndFor
%         \State \Return $M$
%     \end{algorithmic}
% \end{algorithm}
%
% After the pseudocode, walk through the key steps in prose and
% state time/space complexity if it matters to the claims.

\subsection{Implementation Details}
\label{subsec:implementation}

% Hyperparameters, libraries, hardware, engineering choices that
% affect reproducibility — and the REASON for each non-obvious
% choice. Full configs belong in an artifact/appendix, not here.
